\chapter*{Avant-propos}
\addcontentsline{toc}{chapter}{Avant-propos}
\markboth{Avant-propos}{}

\begin{quote}
\itshape « Un livre de mathématiques ne devrait effrayer personne. »
\end{quote}

L'intelligence artificielle repose entièrement sur des mathématiques. Le
problème, c'est que les livres qui les enseignent supposent presque tous
que vous les connaissez déjà : ils vont vite, ils écrivent en symboles
avant d'avoir expliqué en mots, et ils supposent une aisance en
programmation que peu de lecteurs ont au départ. Résultat : beaucoup de
gens curieux, capables, motivés, s'arrêtent à la première page.

\textbf{Ici, on ne suppose presque rien.} Chaque notion est d'abord
racontée avec des mots simples, puis illustrée par une image ou un
dessin, puis seulement écrite en langage mathématique — et enfin mise en
pratique avec quelques lignes de code que vous pouvez exécuter vous-même.

\begin{enclair}
Notre promesse : si vous savez additionner, multiplier et lire un tableau
de nombres, vous pouvez comprendre les mathématiques de l'intelligence
artificielle. Le reste, on le construit ensemble, dans l'ordre, sans
sauter d'étape.
\end{enclair}

\section*{Pour qui est ce livre}

Pour toute personne qui veut appliquer les mathématiques ailleurs —
informatique, IA, traitement du signal, analyse de données — et qui part
d'un niveau nul, ancien ou fragile. Aucun diplôme n'est demandé. Le texte
est écrit pour rester lisible par un élève de début de collège lorsque le
sujet le permet, et pour ne jamais perdre le lecteur lorsque le sujet
devient réellement difficile.

C'est aussi un \emph{tout-en-un} : plutôt que d'aller chercher l'algèbre
linéaire dans un livre, les probabilités dans un deuxième et
l'optimisation dans un troisième — chacun avec ses notations et ses
habitudes —, vous trouverez ici l'ensemble, avec des notations
cohérentes d'un bout à l'autre.

\section*{Comment lire ce livre}

Le texte utilise quelques encadrés récurrents, faciles à reconnaître à
leur pictogramme :

\begin{itemize}
  \item \textbf{\faIcon{comments}~En clair} — l'idée expliquée avec des
        mots de tous les jours.
  \item \textbf{\faIcon{brain}~L'intuition} — l'image mentale à garder en
        tête.
  \item \textbf{\faIcon{robot}~À quoi ça sert en IA ?} — où cette notion
        apparaît vraiment dans un algorithme d'apprentissage.
  \item \textbf{\faIcon{chart-line}~Sur des données réelles} — la même
        idée, appliquée à un jeu de données concret.
  \item \textbf{\faIcon{seedling}~Ce qu'il faut savoir avant} — les seuls
        prérequis de la section, et où les réviser.
  \item \textbf{\faIcon{exclamation-triangle}~Piège à éviter} — l'erreur
        classique que tout le monde fait.
  \item \textbf{\faIcon{key}~À retenir} — l'essentiel, si vous ne deviez
        garder qu'une chose.
  \item \textbf{\faIcon{rocket}~Pour aller plus loin} — un pont vers des
        sujets plus avancés, que vous pouvez ignorer en première lecture.
\end{itemize}

Les résultats importants sont énoncés comme \emph{théorèmes}, et le code
qui les met en œuvre (surtout en Python) est fourni et commenté. Les
démonstrations sont là pour ceux qui veulent savoir \emph{pourquoi} :
sautez-les sans remords à la première lecture, le fil du texte reste
compréhensible sans elles.

\vspace{1em}
\hfill\textit{Bonne lecture.}
\clearpage
